\chapter{One-factor diffusion models}
All one-factor models assume that the short rate is the only state variable. That is, $r_t$ contains
all relevant information about the term structure. Assume the short rate follows a diffusion
process
\[
\d  r_{t}=\alpha\left(r_{t}\right) \d  t+\beta\left(r_{t}\right) \d  z_{t}
\]
where $\left(z_{t}\right)_{t \geq 0}$ is a one-dimensional standard Brownian motion (under the
real-world measure $\P$). The short rate dynamics under the risk-neutral measure
$\Q$ is
\[
\d  r_{t}=\hat{\alpha}\left(r_{t}\right) \d  t+\beta\left(r_{t}\right) \d  z_{t}^{\Q}
\]
where $\hat{\alpha}\left(r_{t}\right)=\alpha\left(r_{t}\right)-\beta\left(r_{t}\right)
\lambda\left(r_{t}\right)$, and where $\lambda$ is the market price of risk (risk premium per unit
volatility). 

Note that these models are called \emph{time homogeneous} since $\hat{\alpha}, \beta$ only depend on
time through $r_{t}$ and not $t$.

\section{Affine models}
A model is affine if the risk-neutral drift rate and the variance rate are affine functions of $r$, i.e.
\begin{equation}
\hat{\alpha}\left(r_{t}\right)=\hat{\varphi}-\hat{\kappa} r, \quad \text { and } \quad \beta(r)^{2}=\delta_{1}+\delta_{2} r \geq 0,
\label{eq:02-affine}
\end{equation}
where $\hat{\varphi}, \hat{\kappa}, \delta_{1}, \delta_{2}$ are constants. When $\hat{\alpha}$ and
$\beta^{2}$ are linear in $r$ we can split the PDE into two ODEs.

\section{Main proof}
\begin{theorem*}[7.1 (Zero-Coupon Bond Prices)]
In an affine model, zero-coupon bond prices are given by
\[
B_{t}^{T}=B^{T}(r, t)=\e^{-a(T-t)-b(T-t) r}
,\]
where the functions $a(\tau)$ and $b(\tau)$ solve $a(0)=b(0)=0$ and
\begin{align}
\frac{1}{2} \delta_{2} b(\tau)^{2}+\hat{\kappa} b(\tau)+b^{\prime}(\tau)-1 & =0  
\label{eq:02-ricatti}
\\
a^{\prime}(\tau)-\hat{\varphi} b(\tau)+\frac{1}{2} \delta_{1} b(\tau)^{2} & =0, \quad \tau>0
\notag
\end{align}
\emph{These ODEs are known as Ricatti equations. For certain parameters $a$ and $b$ can be found
explicitly.}
\end{theorem*}
\begin{proof}
 We prove that the candidate for a solution, is a solution and then use the uniqueness
 of a solution for a PDE. Thus, we show that $B^{T}(r, t)$ solves the fundamental PDE ($\mathcal{S}$
 is state space of $r$)
\begin{align}
\frac{\partial B^{T}(r, t)}{\partial t}
+\hat{\alpha}(r) \frac{\partial B^{T}(r, t)}{\partial r}
+\frac{1}{2} \beta(r)^{2} \frac{\partial^{2} B^{T}(r, t)}{\partial r^{2}}-r B^{T}(r, t) & =0,  
\label{eq:02-fundamental-pde}
\\
B^{T}(r, T) & =1, \quad(r, t) \in \mathcal{S} \times[0, T] \notag
.\end{align}
The terminal condition is obviously satisfied since
\[
B^{T}(r, T)=\e^{-a(0)-b(0) r}=1
\]
Next, we find the derivatives
\begin{align*}
\frac{\partial B^{T}(r, t)}{\partial t} & =B^{T}(r, t)\left(a^{\prime}(T-t)+b^{\prime}(T-t) r\right), \\
\frac{\partial B^{T}(r, t)}{\partial r} & =-B^{T}(r, t) b(T-t), \\
\frac{\partial^{2} B^{T}(r, t)}{\partial r^{2}} & =B^{T}(r, t) b(T-t)^{2} .
\end{align*}
Inserting in \Cref{eq:02-fundamental-pde} and dividing by $B^{T}(r, t)$ we obtain
\[
a^{\prime}(T-t)+b^{\prime}(T-t) r-\hat{\alpha}(r) b(T-t)+\frac{1}{2} \beta(r)^{2} b(T-t)^{2}-r=0 .
\]
Using \Cref{eq:02-affine} and gathering terms we find
\begin{align*}
& a^{\prime}(T-t)+b^{\prime}(T-t) r-(\hat{\varphi}-\hat{\kappa} r) b(T-t)
+\frac{1}{2}\left(\delta_{1}+\delta_{2} r\right) b(T-t)^{2}-r=0 \\
& \implies \quad  
a^{\prime}(T-t)-\hat{\varphi} b(T-t)+\frac{1}{2} \delta_{1} b(T-t)^{2}
+\left(b^{\prime}(T-t)+\hat{\kappa} b(T-t)+\frac{1}{2} \delta_{2} b(T-t)^{2}-1\right)
r=0,
\end{align*}
for all $(r, t) \in \mathcal{S} \times[0, T]$.
Inserting $\tau=T-t$ we observe that the above holds when \Cref{eq:02-ricatti} holds.

\qed
\end{proof}

\section{Vasicek's model}
In Varsicek's model we assume that the short rate $r_{t}$ is driven by an Ornstein-Uhlenbeck process
under $\P$:
\[
\d r_{t}=\kappa\left(\theta-r_{t}\right) \d t+\beta \d z_{t}
\]
where $\kappa, \theta$ and $\beta$ are positive constants. Given the market price of risk
$\lambda(r, t)$, the $\Q$-dynamics are
\[
\d  r_{t}=\left[\kappa\left(\theta-r_{t}\right)-\beta \lambda(r, t)\right] \d t+\beta
\d  z_{t}^{\Q}=\kappa\left(\hat{\theta}-r_{t}\right) \d t+\beta \d 
z_{t}^{\Q}
\]
Assuming $\lambda(r, t)=\lambda$, i.e. constant market price of risk, we have
$\hat{\theta}=\theta-\frac{\beta \lambda}{\kappa}$.

\subsection{Properties}
$r$ is also an Ornstein–Uhlenbeck process under $\Q$, but in this case it is mean-reverting around
$\hat{\theta}$. The model is Gaussian, implying that $r_T$ can take negative values with positive
probability, which is often regarded as a drawback of the model.

The model furthermore assumes constant volatility. It is an affine term structure model with
parameters
\[
\hat{\varphi}=\kappa \hat{\theta}, \quad \hat{\kappa}=\kappa, \quad
\delta_1=\beta^2, \qquad \delta_2=0,
\]
and therefore the price of a zero-coupon bond is given by
\[
B^{T}(r, t)=\e^{-a(T-t)-b(T-t) r},
\]
where the equations in \Cref{eq:02-ricatti} can be solved to:
\[
b(\tau)=\frac{1}{\kappa}\left(1-\e^{-\kappa \tau}\right) \quad \text { and } \quad
a(\tau)=\left(\hat{\theta}-\frac{\beta^{2}}{2 \kappa^{2}}\right)(\tau-b(\tau))
+\frac{\beta^{2}}{4\kappa} b(\tau)^{2}
\]

\subsection{Relevant results}
\begin{description}
\item[Zero-Coupon yields] 
\end{description}
For affine and Vasicek the yield curve satisfies
\[
y^{T}(r, t)=-\frac{\ln \left(B^{T}(r, t)\right)}{T-t}=\frac{a(T-t)}{T-t}+\frac{b(T-t)}{T-t} r .
\]
Affine function of the short rate. For $b(\tau)>0$ yields are increasing in the short rate, an
increase in $r$ shifts the curve upwards (unless $b(\tau)$ proportional to $\tau$ ), prices
decreasing in the short rate. Not all shapes are possible, cannot generate trough and twists.

The instantaneous forward rates (at time $T$) for affine model (maybe Vasicek) is
\[
f^{T}(r, t)=-\frac{\frac{\partial B^{T}(r, t)}{\partial T}}{B^{T}(r,
t)}=a^{\prime}(T-t)+b^{\prime}(T-t) r
\]
\begin{description}
\item[Cox-Ingersoll-Ross (CIR) model] 
\end{description}
Market price of risk is $\frac{\lambda \sqrt{r}}{\beta}$, $\Q$-dynamics is square root proces
\[
\d  r_{t}=\hat{\kappa}\left(\hat{\theta}-r_{t}\right) \d  t+\beta \sqrt{r_{t}} \d  z_{t}^{\Q}
\]
Mean reverting (around $\hat{\theta}$). Non-negative interest rates, 0 is reflecting barrier. Affine
with $\hat{\varphi}=\hat{\kappa} \hat{\theta}, \hat{\kappa}=\hat{\kappa}, \delta_{1}=0$ and
$\delta_{2}=\beta^{2}$ (ODEs can be solved). Time varying volatility (increases with level).


\begin{description}
\item[Generalized affine models] 
\end{description}
In Vasicek we assume the market price of risk is constant. Can let it be affine function of $r$ to
allow for more flexibility. Still affine. Same way of thinking can be used for CIR.

\begin{description}
\item[Non Affine models] 
\end{description}

\begin{itemize}
  \item Quadratic models: Dynamics $r_{t}=x_{t}^{2}$ ($x_{t}$ Ornstein-Uhlenbeck process). $a$ and $b$ as in affine case. 0 is refelcting barrier (disadvantage). $B^{T}(x, t)= \e^{-a(T-t)-b(T-t) x-c(T-t) x^{2}}$. Complicated closed form expressions for prices of derivatives.
  \item Shadow-rate models: Dynamics $r_{t}=\max \left\{0, x_{t}\right\}, x_{t}$ OU process. Only nonnegative values. 0 is absorbing barrier (occurs in practice). Bond and derivative prices computed numerically.
\end{itemize}
